\documentclass{article}
\usepackage[UTF8]{ctex}
\usepackage{amsmath,amssymb,amsfonts}
\usepackage{enumitem}
\usepackage{geometry}
\geometry{a4paper, margin=1in}
% === 错题本排版优化（不参与题目标识） ===
\setlength{\parindent}{0pt}
\setlength{\parskip}{0.35em plus 0.10em minus 0.08em}
\setlength{\emergencystretch}{3em}
\setlist{itemsep=0.25em, topsep=0.40em, parsep=0pt}
\allowdisplaybreaks[2]
% === 错题本排版优化结束 ===
\begin{document}
\section*{东华大学2024年高等代数真题}
\section*{41．东华大学2024年高等代数试卷}

\subsection{题目1}
1．（15 分）设 \(A, B\) 分别为 \(m \times n\) 与 \(n \times m\) 矩阵，记 \(M=\left(\begin{array}{cc}E_{m} & A \\ B & E_{n}\end{array}\right)\) ．

\begin{itemize}
  \item[（1）] 证明：\(\left|E_{m}-B A\right|=\left|E_{n}-A B\right|\) ．
  \item[（2）] 记 \(N=\left(\begin{array}{cc}\lambda E_{m} & A \\ B & \lambda E_{n}\end{array}\right)\) ，通过求 \(|N|\) 讨论 \(A B\) 与 \(B A\) 特征值的关系．
\end{itemize}

\subsection{题目2}
2．（20 分）设矩阵

\[
A=\left(\begin{array}{lllll}
1 & & & & 1 \\
1 & 1 & & & \\
& 1 & \ddots & & \\
& & \ddots & 1 & \\
& & & 1 & 1
\end{array}\right) .
\]

求 \(A\) 的行列式，秩和 \(A^{*}\) 。

\subsection{题目3}
3．（20 分）设 \(A\) 为可逆实矩阵。

\begin{itemize}
  \item[（1）] 证明：\(A^{\mathrm{T}} A\) 为正定矩阵。
  \item[（2）] 证明：\(A=C D\) ，其中 \(C\) 是正交矩阵，\(D\) 为正定矩阵。
\end{itemize}

\subsection{题目4}
4．（20 分）设 \(V\) 是数域 \(P\) 上所有二阶矩阵构成的线性空间，\(B=\left(\begin{array}{cc}1 & -1 \\ 2 & 1\end{array}\right) \in V\) ，定义 \(V\) 上的线性变换 \(\sigma: \sigma(X)=B^{-1} X B\) ，求 \(\sigma\) 的所有特征子空间。

\subsection{题目5}
5．（15 分）设 \(V\) 为全体 \(n\) 阶实对称矩阵构成的线性空间，定义 \((A, B)=\operatorname{tr}(A B), A, B \in V\) 。

\begin{itemize}
  \item[（1）] 证明：\(V\) 满足欧氏空间的条件．
  \item[（2）] 记子空间 \(W=\{X \in V \mid \operatorname{tr}(X)=0\}\) ，求 \(W\) 的基与维数及正交补空间 \(W^{\perp}\) 的基与维数。
\end{itemize}

\subsection{题目6}
6．（25 分）对于线性空间 \(V\) 上的线性变换 \(\sigma\) ，若存在正整数 \(m\) ，使得 \(\sigma^{m}=0\) ，则称 \(\sigma\) 为幂零变换．

\begin{itemize}
  \item[（1）] 证明：\(\sigma\) 是幂零变换当且仅当 \(\sigma\) 的特征值全为 0 。
  \item[（2）] 若 \(\sigma, \tau\) 均为幂零变换，且 \(\sigma \tau=\tau \sigma\) ，证明 \(\sigma+\tau\) 为幂零变换．
  \item[（3）] \(\sigma\) 为幂零变换，\(\tau\) 为可逆变换，且 \(\sigma \tau=\tau \sigma\) ，证明 \(\sigma+\tau\) 为可逆变换．
\end{itemize}

\subsection{题目7}
7．（20 分）设 \(f\) 为 \(V \rightarrow Q\) 上的变换，其中 \(V\) 为有理数域上 \(n\) 阶矩阵构成的线性空间。

\begin{itemize}
  \item[（1）] 若 \(f(A+B)=f(A)+f(B), \forall A, B \in V\) ，证明：\(f\) 为线性函数．
  \item[（2）] 若 \(f\) 满足（1）的条件，且 \(f(A B)=f(B A), \forall A, B \in V\) ，证明：
\[
f(X)=c \sum_{i=1}^{n} a_{i i}
\]
其中 \(X=\left(a_{i j}\right)_{n \times n} \in V, c\) 为有理数．
\end{itemize}

\subsection{题目8}
8．（15 分）已知 \(f, g\) 是线性空间 \(V\) 上的线性变换，证明：

\begin{itemize}
  \item[（1）] \(\operatorname{Im} f \subseteq \operatorname{Im} g\) 当且仅当存在线性变换 \(h\) 使得 \(g=f h\) 。
  \item[（2）] \(\operatorname{Ker} f \subseteq \operatorname{Ker} g\) 当且仅当存在线性变换 \(h\) 使得 \(f=g h\) ．
\end{itemize}
\end{document}
